\chapter{Dos, y no uno solo}

\epigraph{Lo contrario del amor no es el odio, sino la indiferencia. Y la indiferencia solo es posible ante lo que no me interpela, ante lo que es igual a mí.}{Adaptado de Elie Wiesel}

Pocos temas generan tanta tensión hoy como la diferencia entre el varón y la mujer. Y se entiende. Durante siglos, esa diferencia se utilizó para justificar la subordinación de la mujer: si somos ``complementarios'', decía el argumento, cada uno tiene su lugar ---y el lugar de ella resultaba ser siempre el segundo---. Cualquier persona honesta tiene que reconocer que eso ocurrió, que hizo mucho daño, y que la reacción contra esa injusticia ha sido en gran parte legítima.

Pero hay una pregunta que merece la pena hacerse con calma, sin trincheras y sin eslóganes: ¿y si el problema no fuera la diferencia, sino lo que hicimos con ella? ¿Y si la complementariedad, bien entendida, no tuviera nada que ver con la desigualdad, sino con algo mucho más hondo ---con la posibilidad misma del encuentro---?

Este capítulo no pretende zanjar un debate que sigue abierto. Pretende algo más modesto: mirar la diferencia sexual con ojos limpios. No como un campo de batalla, sino como un lugar de asombro.

\section{Lo que la complementariedad no es}

Conviene empezar por despejar el terreno. Cuando aquí hablamos de complementariedad no estamos diciendo que el varón sea la pieza A y la mujer la pieza B de un puzzle donde cada uno tiene una función predeterminada. No estamos diciendo que ella sea la emoción y él la razón, que ella cuide y él provea, que ella sienta y él decida. Esos moldes rígidos han causado mucho sufrimiento y, además, son falsos: cualquiera que conozca a hombres y mujeres reales sabe que la realidad es infinitamente más rica que el estereotipo.

Tampoco estamos diciendo que la diferencia sexual sea un sistema jerárquico donde uno vale más que el otro. La igual dignidad del varón y la mujer no es una concesión moderna ni un gesto de cortesía: es un dato antropológico fundamental. Somos igualmente humanos, igualmente dignos, igualmente capaces de razón, de libertad y de amor.\footnote{La igualdad radical en dignidad entre varón y mujer es afirmada hoy tanto desde la antropología filosófica como desde las ciencias sociales. Desde una perspectiva creyente, el Génesis lo expresa con nitidez: ``varón y mujer los creó'' (Gn 1,27), ambos imagen de Dios, sin jerarquía alguna en el acto creador.}

Lo que sí estamos diciendo es algo diferente: que el varón y la mujer, siendo iguales en dignidad, no son idénticos. Y que esa no-identidad, lejos de ser un problema, es una riqueza. Quizá la mayor riqueza de la condición humana.

\section{Un dato que precede a la cultura}

Vivimos en un tiempo que tiende a pensar que toda diferencia entre hombres y mujeres es cultural, construida, y por tanto modificable a voluntad. Es una postura comprensible: si las diferencias son solo culturales, eliminar la desigualdad es cuestión de cambiar la cultura. Y algo de razón tiene: muchas de las diferencias que durante siglos se consideraron ``naturales'' ---que las mujeres no podían estudiar, votar, o dirigir--- eran efectivamente culturales y, con justicia, han sido superadas.

Pero la ciencia nos invita a no simplificar demasiado. Porque junto a las diferencias que son culturales hay otras que no lo son ---o al menos, no del todo---.

Empecemos por lo más básico. Cada célula del cuerpo de un varón tiene un par de cromosomas XY; cada célula de una mujer, XX. Esa diferencia genética no se limita a los órganos reproductivos: influye en la expresión de miles de genes a lo largo de todo el organismo.\footnote{Cf. Arnold, A. P., ``A general theory of sexual differentiation'', \textit{Journal of Neuroscience Research}, 2017. Arnold argumenta que los efectos de los cromosomas sexuales van mucho más allá de las gónadas y afectan a prácticamente todos los tejidos del cuerpo.} Durante el desarrollo embrionario, la exposición a diferentes niveles de hormonas ---especialmente testosterona y estrógenos--- organiza el cerebro de formas que, sin ser determinantes, son significativas.

Los estudios de neurociencia muestran diferencias estadísticas entre cerebros masculinos y femeninos: en el tamaño relativo de ciertas estructuras, en la conectividad entre hemisferios, en la densidad de receptores de determinados neurotransmisores.\footnote{El estudio más citado es el de Ingalhalikar et al., ``Sex differences in the structural connectome of the human brain'', \textit{PNAS}, 2014, que encontró patrones de conectividad distintos entre cerebros masculinos y femeninos. Es un hallazgo sugerente pero debatido: otros investigadores han señalado que los tamaños de efecto son pequeños y que el solapamiento entre individuos de ambos sexos es mayor que la diferencia entre las medias. Para una revisión más amplia y matizada, cf. Halpern, D.\,F., \textit{Sex Differences in Cognitive Abilities}, Psychology Press, 2012, donde se concluye que las diferencias cognitivas entre sexos son reales pero modestas, parcialmente biológicas y parcialmente moldeadas por la experiencia.} El perfil hormonal es diferente a lo largo de toda la vida: los niveles de testosterona, estrógenos, progesterona y oxitocina varían notablemente entre varones y mujeres, y esas hormonas influyen en la percepción del dolor, en la respuesta al estrés, en los patrones de sueño, en la forma de procesar las emociones.\footnote{Cf. Taylor et al., ``Biobehavioral responses to stress in females: Tend-and-befriend, not fight-or-flight'', \textit{Psychological Review}, 2000. Este influyente artículo mostró que la respuesta al estrés en mujeres tiende a seguir un patrón de ``cuidar y buscar alianzas'' (\textit{tend and befriend}), mediado por la oxitocina, a diferencia del clásico modelo masculino de ``luchar o huir''.}

Nada de esto significa que los hombres sean ``de Marte'' y las mujeres ``de Venus''. La variabilidad dentro de cada sexo es enorme, y la cultura moldea poderosamente la expresión de estas tendencias. Un hombre puede ser profundamente empático y una mujer enormemente analítica, sin que ninguno de los dos sea menos hombre o menos mujer por ello. Las diferencias biológicas son tendencias estadísticas, no destinos individuales.

Pero sería igualmente ingenuo negar que esas tendencias existen. Y lo más interesante no es que existan, sino para qué existen.

Porque la diferencia entre varón y mujer no se agota en los cromosomas, las hormonas o las conexiones neuronales. Todo eso es real, pero es el signo, no el significado. Lo que esos datos biológicos están señalando es algo más profundo: que ser varón y ser mujer son dos maneras distintas de ser persona. No solo dos cuerpos diferentes, sino dos formas de estar en el mundo, de percibir, de abrirse a la realidad. Y cuando esas dos formas se encuentran de verdad, ocurre algo que la soledad ---por rica que sea--- no puede dar: cada uno descubre en el otro dimensiones de lo humano a las que, desde sí mismo, no tenía acceso. La diferencia no es un accidente biológico que hay que tolerar; es la condición que hace posible el regalo más grande que una persona puede hacer: entregarse a alguien que es genuinamente otro, y recibir de ese otro lo que uno no podía darse a sí mismo. Sin alteridad real, el amor se convierte en un espejo ---bonito, quizá, pero estéril---. Solo ante alguien verdaderamente distinto puedo salir de mí, arriesgarme, y descubrir que dar la vida tiene sentido.

\section{La alteridad como condición del encuentro}

Aquí es donde la reflexión se pone interesante. Porque si nos preguntamos por qué la naturaleza ---o, si se prefiere, la evolución--- ha mantenido la diferencia sexual en una especie tan sofisticada como la nuestra, la respuesta no es solo ``para la reproducción''. La reproducción sexual existe en organismos que no tienen nada parecido a lo que nosotros llamamos encuentro. Pero en el ser humano, la diferencia sexual hace posible algo más: hace posible que haya un \textit{otro} de verdad.

Pensemos en esto despacio. ¿Qué pasa cuando me encuentro con alguien que es exactamente igual a mí? Que, en el fondo, no me encuentro con nadie. Me encuentro con un espejo. Y un espejo no me interpela, no me desafía, no me saca de mí mismo. Me confirma. Me deja donde estoy.

El verdadero encuentro ---ese que me transforma, que me obliga a salir de mi mundo y entrar en el de otro--- solo es posible cuando hay alteridad, cuando el otro es realmente \textit{otro}. Alguien que ve las cosas de un modo que yo no había previsto. Que siente de una manera que no es la mía. Que me revela dimensiones de la realidad a las que yo solo, desde mi perspectiva, no tengo acceso.

La diferencia sexual es, quizá, la forma más radical y más cotidiana de alteridad que experimentamos los seres humanos. No la única, pero sí la más constitutiva. Cuando un varón y una mujer se encuentran de verdad ---no como estereotipos, sino como personas concretas, con toda su complejidad---, cada uno tiene la oportunidad de descubrir un modo de ser humano que no es el suyo. Y eso es un regalo, no una amenaza.\footnote{Emmanuel Levinas desarrolló filosóficamente esta idea de la alteridad como condición del encuentro ético. Para Levinas, el rostro del otro ---irreductiblemente diferente de mí--- es lo que funda la relación ética y me saca de mi soledad. Cf. Levinas, E., \textit{Totalidad e infinito}, Sígueme, 1977.}

Es importante decir algo aquí con honestidad. Al hablar de la riqueza de la diferencia sexual en el matrimonio, no estamos negando que exista amor verdadero, generoso y sacrificado en otras formas de relación. La experiencia nos muestra personas que aman con autenticidad y entrega a personas de su mismo sexo, y ese amor ---que es libre, que busca el bien del otro, que es capaz de sacrificio--- merece respeto y comprensión.\footnote{El \textit{Catecismo de la Iglesia Católica} pide expresamente que las personas homosexuales «deben ser acogidas con respeto, compasión y delicadeza» y que «se evitará, respecto a ellos, todo signo de discriminación injusta'' (CIC, n. 2358). El Papa Francisco ha insistido en múltiples ocasiones en que la Iglesia no puede ser una aduana sino una casa abierta.} Lo que este capítulo explora es algo específico: la riqueza particular que la diferencia sexual aporta al encuentro conyugal. Del mismo modo, tampoco excluimos otras vocaciones orientadas al amor: la vocación sacerdotal, la vida consagrada, quienes eligen el celibato por el Reino o por otras razones ---todas ellas son caminos auténticos de entrega y de amor---. Este libro habla de una vocación particular, la del matrimonio entre varón y mujer, que tiene una característica singular: es la más común dentro del género humano. La inmensa mayoría de las personas que han vivido, viven y vivirán están llamadas a amar de esta manera concreta. Hablar de ella con profundidad no es menospreciar las demás. Es, sencillamente, mirar con asombro lo que la diferencia entre varón y mujer hace posible, sin que eso signifique que el amor solo exista ahí.

\section{La riqueza que nace de la diferencia}

Bajemos de la filosofía a la vida concreta. Cualquier pareja con algunos años de convivencia sabe de qué estamos hablando, aunque quizá nunca lo haya formulado así.

Ella llega a casa preocupada por un problema en el trabajo. No quiere que él se lo resuelva ---quiere contárselo, necesita que la escuche, que valide lo que siente---. Él, mientras tanto, escucha el problema y automáticamente empieza a buscar soluciones. No porque no le importe lo que ella siente, sino porque así funciona su cabeza: si hay un problema, hay que resolverlo. Los dos tienen razón. Los dos están expresando formas legítimas de afrontar la dificultad. Y los dos necesitan aprender del otro: ella, a veces, necesita que le ayuden a actuar; él, a veces, necesita que le ayuden a parar y simplemente estar.

¿Es este ejemplo un estereotipo? Parcialmente, sí. No todas las mujeres procesan así ni todos los hombres son así. Pero la investigación en psicología de la comunicación muestra que estas tendencias existen como patrones estadísticos, y que buena parte de los conflictos de pareja nacen precisamente de no entender que el otro no está fallando ---está siendo diferente---.\footnote{Cf. Tannen, D., \textit{You Just Don't Understand: Women and Men in Conversation}, William Morrow, 1990. Tannen, lingüista de la Universidad de Georgetown, documentó extensamente las diferencias en los estilos comunicativos de hombres y mujeres, subrayando que la mayoría de los malentendidos nacen de asumir que el otro debería comunicarse como uno mismo.}

La diferencia obliga a salir de uno mismo. Obliga a escuchar, a intentar comprender una lógica que no es la propia, a aceptar que mi forma de ver el mundo no es la única forma de verlo. Y eso, que a veces resulta agotador, es exactamente lo que nos hace crecer.

Una pareja donde los dos son idénticos ---donde piensan igual, sienten igual, reaccionan igual--- sería una pareja cómoda, quizá, pero estéril. No en el sentido biológico, sino en el sentido humano: no habría fricción creativa, no habría ese choque fecundo entre dos maneras de ver que obliga a construir algo nuevo, algo que ninguno de los dos habría podido construir solo.

\section{El cuerpo que habla}

Hay algo más. La diferencia sexual no es solo una diferencia psicológica o de carácter. Es, ante todo, una diferencia inscrita en el cuerpo. Y el cuerpo humano, como hemos visto en los capítulos anteriores, no es un accidente: es un lenguaje.

Nuestro cuerpo habla. Habla cuando una madre abraza a su hijo y la oxitocina inunda a los dos. Habla cuando el llanto de un bebé activa circuitos de alarma en el cerebro de sus padres. Habla cuando el contacto piel con piel entre dos personas que se quieren baja el cortisol y sube la sensación de seguridad.

Y habla también a través de la diferencia sexual. El cuerpo del varón y el de la mujer ``dicen'' cosas diferentes. No mejores ni peores: diferentes. La configuración hormonal de ella la inclina ---en términos estadísticos, insisto--- hacia una mayor sensibilidad a las señales emocionales del otro, hacia una mayor facilidad para el lenguaje verbal de las emociones, hacia una respuesta al estrés que busca la conexión.\footnote{Cf. la ya mencionada hipótesis \textit{tend and befriend} de Taylor et al. Estudios posteriores han confirmado que el sistema oxitocina-estrógenos favorece respuestas prosociales ante el estrés, especialmente en mujeres.} La configuración hormonal de él lo inclina hacia una mayor orientación espacial, hacia una respuesta al estrés más orientada a la acción, hacia una protectividad que se expresa muchas veces más con hechos que con palabras.\footnote{Cf. Archer, J., ``The reality and evolutionary significance of human psychological sex differences'', \textit{Biological Reviews}, 2019. Archer ofrece una revisión equilibrada de las diferencias psicológicas entre sexos, distinguiendo las que tienen base evolutiva sólida de las que son más culturales.}

Estas inclinaciones no son cárceles. Son puntos de partida. Cada persona concreta las modula, las amplía, las matiza con su historia, su educación, su libertad. Pero negar que existen es negar lo que el cuerpo dice. Y en un libro que ha empezado escuchando lo que el cuerpo tiene que decirnos sobre el amor, sería incoherente dejar de escucharlo aquí.

\section{Ni uniformidad ni guerra}

Si tuviéramos que resumir lo que hemos visto en este capítulo, podríamos decirlo así: la diferencia sexual humana se mueve entre dos peligros. Uno es antiguo: usar la diferencia para justificar la desigualdad, asignando roles rígidos y jerarquías donde no las hay. El otro es más reciente: negar la diferencia en nombre de la igualdad, como si reconocer que varones y mujeres no son idénticos fuera automáticamente una forma de discriminación.\footnote{Cf. Badinter, E., \textit{XY: De l'identité masculine}, Odile Jacob, 1992; y en una línea diferente, Halpern, D. F., \textit{Sex Differences in Cognitive Abilities}, Psychology Press, 2012. Ambas autoras, desde perspectivas distintas, ilustran la tensión entre reconocer las diferencias y evitar que se conviertan en herramientas de opresión.}

La realidad es más sutil que ambos extremos. Somos iguales en dignidad y diferentes en el modo de ser. Y esa diferencia no es un error a corregir ni una jerarquía a imponer: es una invitación al encuentro.

Pensemos en una orquesta. Que el violín y el violonchelo sean instrumentos diferentes no significa que uno sea superior al otro. Pero tampoco significa que den la misma nota. La belleza de la música nace precisamente de que cada uno aporta un timbre, un registro, una textura que el otro no tiene. Si todos los instrumentos fueran iguales, no habría armonía: habría unísono. Y el unísono, por potente que sea, es infinitamente más pobre que la sinfonía.

El varón y la mujer son como dos instrumentos distintos que, cuando afinan juntos, producen algo que ninguno de los dos podría producir solo. Eso es la complementariedad bien entendida. No la media naranja que me completa porque estoy incompleto, sino la otra voz que enriquece la canción precisamente porque suena diferente.

\section{Ante el misterio del otro}

Hay una última cosa que conviene decir, y quizá sea la más importante.

La diferencia sexual nos sitúa, permanentemente, ante el misterio del otro. Porque por mucho que conviva con mi pareja, por mucho que la conozca, por mucho que hayamos compartido, siempre habrá algo en ella ---o en él--- que se me escapa. Una forma de sentir que no es la mía. Una manera de estar en el mundo que no puedo experimentar desde dentro. Un misterio que no se agota.

Y eso, lejos de ser una frustración, es la garantía de que el amor no se aburre. Porque lo que se posee del todo deja de interesar. Lo que se domina completamente deja de fascinar. Solo lo que conserva algo de misterio sigue invitándome a acercarme, a descubrir, a asombrarme.

La diferencia sexual nos recuerda que el otro no es una extensión de mí mismo. No es mi reflejo, ni mi complemento mecánico, ni la pieza que me faltaba. Es alguien. Alguien irreductiblemente distinto, irreductiblemente libre, irreductiblemente otro. Y es justamente por eso que puedo amarlo de verdad. Porque el amor, el verdadero, no es fusión: es encuentro. Y para que haya encuentro, tiene que haber dos.

Dos, y no uno solo.

\paracaminar

\begin{enumerate}
  \item Pensad juntos: ¿en qué momentos concretos de vuestra relación habéis experimentado la diferencia del otro como una riqueza? ¿Y en qué momentos os ha costado más aceptarla?

  \item ¿Hay algo de vuestra pareja que, por más que lo intentéis, no termináis de entender del todo? ¿Podéis ver ese ``misterio'' no como un problema, sino como una invitación a seguir conociéndoos?

  \item Hablad con sinceridad: ¿alguna vez habéis caído en la tentación de querer que el otro piense, sienta o reaccione como vosotros? ¿Qué pasaría si, en lugar de eso, agradeciérais la diferencia?
\end{enumerate}

\vinetacapitulo{ch05}
